\begin{proof}[Proof of \Cref{lem:p031-rigsmooth-curve-envelope}]
Put \(s:=\sup_i u_i\) and \(u:=s^*\).  Local boundedness from above implies
that \(u\) takes values in \(\mathbf R\cup\{-\infty\}\).  Upper
semicontinuous regularization is local on open subsets: if \(W\subset C\)
is open, then
\begin{equation}
\label{eq:p031-usc-open-restriction}
                         u|_W=(s|_W)^*.
\end{equation}
Indeed, the neighborhoods of a point \(x\in W\) that are contained in
\(W\) are cofinal in the neighborhoods of \(x\) in \(C\), and both sides
of \eqref{eq:p031-usc-open-restriction} equal the corresponding
\(\limsup_{y\to x}s(y)\).

We first work on a connected component \(W\) of
\(\operatorname{Int}(C)\).  It is a smooth, boundaryless, strict analytic
curve.  By \cite[\S4.1.2 and \S3.4.1]{GR25Classical}, every \(u_i|_W\)
is either subharmonic in Thuillier's sense or identically \(-\infty\) on
\(W\).  If \(u|_W\equiv-\infty\), it is subharmonic in the extended sense
of \cite[Definition~3.4]{GR25Classical}.  We may therefore suppose that it
is not identically \(-\infty\).

Let \(D\subset W\) be a strict affinoid domain, and let
\(q:D\to\mathbf R\) be harmonic with
\begin{equation}
\label{eq:p031-harmonic-boundary-majorant}
                              u\leq q
                              \quad\text{on }\partial D.
\end{equation}
Since \(u_i\leq s\leq u\), every \(u_i\) satisfies the same boundary
inequality.  Thuillier's defining comparison property gives
\[
                              u_i\leq q
                              \quad\text{on }D
\]
for every \(i\); the assertion is tautological when \(u_i\equiv-\infty\).
Hence \(s\leq q\) on \(D\).  The complement of the relative boundary of
\(D\) is its topological interior in \(W\): this is the combination of
the relative-interior identities in \cite[\S2.11(3),(4)]{GR25Classical},
using that \(W\) is boundaryless and separated.  Therefore, for
\(x\in D\setminus\partial D\), upper semicontinuous regularization and
continuity of \(q\) give
\[
                         u(x)=\limsup_{y\to x}s(y)
                              \leq q(x).
\]
On \(\partial D\) the same inequality is
\eqref{eq:p031-harmonic-boundary-majorant}.  Thus \(u\leq q\) on \(D\).
Thuillier's comparison criterion proves that \(u|_W\) is subharmonic.

We have proved that \(u\) is subharmonic on
\(\operatorname{Int}(C)\).  The boundary-extension property
\cite[Proposition~3.5(3)]{GR25Classical} makes \(u\) subharmonic on all of
\(C\), including its relative boundary.  Finally,
\cite[\S4.1.2]{GR25Classical} identifies subharmonicity and classical
plurisubharmonicity on analytic curves.  This proves the lemma.
\end{proof}

\begin{proof}[Proof of
\Cref{thm:p031-perfected-unibranch-curve-envelope}]
Write
\[
                  v:=\sup_{i\in I}\varphi_i,
                  \qquad u:=v^*.
\]
We first record a topological identity.  If \(a:X\to
\mathbf R\cup\{-\infty\}\) is locally bounded above and
\(q:Y\to X\) is a homeomorphism, then
\begin{equation}
\label{eq:p031-homeomorphism-usc}
                          a^*\circ q=(a\circ q)^*.
\end{equation}
This follows directly because \(q\) identifies the two neighborhood
filters in the \(\limsup\) definition of upper semicontinuous
regularization.  The identity also commutes with a pointwise supremum:
\((\sup_i\varphi_i)\circ q=\sup_i(\varphi_i\circ q)\).

Let
\[
 \pi:U^\dagger\longrightarrow U,
 \qquad
 \iota:Z\lhook\joinrel\longrightarrow U^\dagger
\]
be the scalar-extension projection and the reduction inclusion.  In
positive characteristic, \(\pi\) is a homeomorphism by
\cite[Lemma~3.7]{GR25Classical}; in characteristic zero it is the identity.
The map \(\iota\) is a homeomorphism by
\cite[Lemma~3.8]{GR25Classical}.  The normalization map \(\nu\) is a
topological quotient by \cite[Lemma~3.9(4)]{GR25Classical}; it is therefore
a homeomorphism under the bijectivity hypothesis of the theorem.
Repeated use of \eqref{eq:p031-homeomorphism-usc} yields
\begin{equation}
\label{eq:p031-upstairs-envelope-identity}
 u\circ\pi\circ\iota\circ\nu
 =
 \left(
   \sup_{i\in I}\varphi_i\circ\pi\circ\iota\circ\nu
 \right)^*.
\end{equation}

The field \(k^\dagger\) is perfect.  Since \(Z\) is a reduced strict
\(k^\dagger\)-analytic curve, its normalization \(\widetilde Z\) is
rig-smooth by \cite[Lemma~3.9(1)]{GR25Classical}.  Each function
\(\varphi_i\circ\pi\circ\iota\circ\nu\) is classically
plurisubharmonic: scalar extension is
\cite[Proposition~4.2(5)]{GR25Classical}, and analytic pullback is the
forward implication of \cite[Proposition~4.2(6)]{GR25Classical}.  The
family remains locally bounded above because all three underlying maps are
homeomorphisms.  Applying
\Cref{lem:p031-rigsmooth-curve-envelope} to
\eqref{eq:p031-upstairs-envelope-identity} shows that
\(u\circ\pi\circ\iota\circ\nu\) is classically plurisubharmonic on
\(\widetilde Z\).

The map \(\nu\) is finite and surjective.  The map \(\iota\) is a finite
closed immersion and is surjective on underlying points.  The finite
surjective descent implication in
\cite[Proposition~4.2(6)]{GR25Classical}, applied first to \(\nu\) and then
to \(\iota\), gives
\[
                         u\circ\pi
                         \quad\text{classically psh on }U^\dagger.
\]
The converse scalar-extension implication in
\cite[Proposition~4.2(5)]{GR25Classical} now shows that \(u\) is
classically psh on \(U\).  This proves
\eqref{eq:p031-perfected-unibranch-envelope}.

If \(k\) is perfect, then \(k^\dagger=k\) and \(Z=U\).  For a normal
curve the normalization is the identity.  More generally, if the
normalization has one point over every point, then it is bijective, which
is precisely the explicit analytic-unibranch condition used in the final
assertion.
\end{proof}

\begin{proof}[Proof of
\Cref{cor:p031-local-fs-envelope-on-unibranch-curves}]
A locally Fubini--Study function is classically plurisubharmonic by the
logarithm, positive-scaling, finite-maximum, and sheaf properties in
\cite[Proposition~4.2(1)--(3)]{GR25Classical}.  A semipositive real-PL
function is classically plurisubharmonic by
\cite[Theorem~4.3]{GR25Classical}.  Thus every member of
\(\mathcal A_f\) is classically psh.

The upper semicontinuous function \(f\), which never takes the value
\(+\infty\), is locally bounded above.  Hence the family
\(\mathcal A_f\) is locally bounded above.  If it is nonempty, apply
\Cref{thm:p031-perfected-unibranch-curve-envelope}.  If it is empty, the
regularized supremum is the constant function \(-\infty\), which is
classically psh by \cite[Definitions~3.4 and~4.1]{GR25Classical}.
Relative-boundary points were retained in
\Cref{lem:p031-rigsmooth-curve-envelope} through
\cite[Proposition~3.5(3)]{GR25Classical}.
\end{proof}
